\subsection*{Coupled A.11/A.19 root compatibility and terminal-gauge identity (22 October 2026)}
\label{subsec:a11-a19-coupled-gauge-20261022}
Distinguish the two azimuthal corrections $q_1,q_2$ in (A.11)
from the separate pulse-velocity moment corrections in (A.9)/(A.19).
The source orders the stages so the later azimuthal field
$E_c=E_b(1+q_1\beta_1+q_2\beta_2)$ is independent of the
pulse amplitude $A$, while the pulse changes $U$ and its two
velocity-moment corrections are affine in $A$.
The exact common-profile residuals are
\[
R_1=u(q_1+xq_2)-\Delta_I(\eta),\quad
R_2=2v(q_1+zq_2)+w(q_1^2+zq_2^2),
\]
\[
R_3=k_2(q,\eta)A^2+k_1(q,\eta)A+k_0(q,\eta).
\]
The $k_i$ may depend on the genuinely corrected azimuthal profile.
At fixed $\lambda$, $u,v,w,x,z$ are independent of $\eta$.

\begin{proposition}[Exact block-triangular Jacobian]
The Jacobian in unknowns $(q_1,q_2,A)$ is
\[
\begin{pmatrix}
u&ux&0\\
2(v+wq_1)&2z(v+wq_2)&0\\
\partial_{q_1}R_3&\partial_{q_2}R_3&2k_2A+k_1
\end{pmatrix},
\]
with determinant
\begin{equation}\label{eq:joint-A11-A19-jac-20261022}
\boxed{2u[z(v+wq_2)-x(v+wq_1)](2k_2A+k_1).}
\end{equation}
If the small A.11 root has nonsingular Jacobian, $k_2>0$
and $k_0<0$ at a positive A.19 root, then
\[
2k_2A+k_1=\sqrt{k_1^2-4k_2k_0}>0.
\]
Thus the combined root is locally unique and $C^1$ when the
coefficients are $C^1$. If the needed roots are unique for every
$\eta\in[-1,1]$, the local branches paste to a common $C^1$ family.
\end{proposition}
\begin{proof}
Expand the determinant along its third column. For
$k_2>0,k_0<0$, the quadratic has exactly one positive
root, with a strictly positive derivative given by the square
root of its discriminant. Apply the implicit function theorem.
\end{proof}

\begin{proposition}[Exact pressure-neutrality factor]
On the source plateau,
$E_b(y)^2=E_b(y_1)^2e^{-(1+2\lambda)(y-y_1)}$,
with the two disjoint translated bumps centered two units apart.
Therefore
\begin{equation}\label{eq:joint-A11-pressure-factor-20261022}
\boxed{\ \int(E_c^2-E_b^2)\,dy=E_b(y_1)^2R_2(q,\eta).\ }
\end{equation}
Consequently the exact A.11 second equation forces the pressure
increment and its \emph{total} $\eta$ derivative to vanish identically.
If the prescribed unedited reference already satisfies its
terminal pressure normalization, the same corrected profile
inherits $\Pi(\infty,\eta)=0$. This baseline gauge is a distinct
source assumption; A.11 neutrality alone merely preserves it.
\end{proposition}
\begin{proof}
Expand the square. Disjointness removes all mixed bump terms.
The first support contributes the linear and quadratic weights
$v,w$; translation of the second support multiplies both by
$z=e^{-2(1+2\lambda)}$. This yields exactly $R_2$.
Differentiate along the common $C^1$ root and use the source
pressure convention
$\Pi(\infty)-\Pi^{\rm base}(\infty)=\frac12\int(E_c^2-E_b^2)\,dy$.
\end{proof}

\begin{proposition}[Exact total-amplitude sensitivity]
Along a common $C^1$ positive root of $R_3=0$, let a dot
denote the total $\eta$ derivative through $q_1,q_2$.
Then
\[
\boxed{\quad
 A_\eta=-\frac{\dot k_2A^2+\dot k_1A+\dot k_0}
 {2k_2A+k_1}
 =-\frac{\dot k_2A^2+\dot k_1A+\dot k_0}
 {\sqrt{k_1^2-4k_2k_0}}.
\quad}
\]
A source-certified lower bound $2k_2A+k_1\ge m_A>0$
therefore gives
$|A_\eta|\le
(A^2|\dot k_2|+A|\dot k_1|+|\dot k_0|)/m_A$.
If $R_3$ is the genuine normalized terminal $S$ equation,
then $S(\infty,\eta)$ and its \emph{total} derivative vanish
on this root. With the preceding exact A.11 pressure equation
and a certified baseline gauge, all four terminal-residual variables
$\mathcal F,\mathcal F_\eta,\mathcal G,\mathcal G_\eta$ vanish.
The terminal-gauge contribution to $N_s$ is then identically zero;
the separate future-tail, other stress cones and physical PDE
remain to be proved.
\end{proposition}
\begin{proof}
Differentiate $R_3=0$ along the joint branch and divide
by the positive $A$ derivative. Apply the total derivative
to the two identical-zero terminal residual functions.
\end{proof}

\paragraph{Adversarial negative control and release boundary.}
At $q_1=1/10$, $q_2=0$, $u=v=1$, $x=7$, $z=1/8$,
$w=20$, $\Delta_I=1/10$, and $R_3=A^2-1$, the
values $A=1$, $R_1=R_3=0$ coexist with $R_2=2/5\ne0$.
Thus an A.19 root \emph{alone} does not establish
pressure neutrality. Independent symbolic and 10,000 rational
synthetic checks confirm the displayed finite-dimensional
identities; they do not authenticate a single source-admissible
parameter tuple or the true terminal baseline pressure.
The global originating source theorem has forcing.
Appendix B, unforced similarity startups, exterior stress
certification, zero-force physical PDE residual and proposed
unforced Theorem~1.1 all remain OPEN/HOLD.
